\chapter{template}
This is the template of the chapter in split file.

\section{s1}
This is a section.

\subsection{sub1}
This is a subsection.

% equation
\begin{equation}
f_n(\bm{r})=
\begin{cases}
\frac{l_n}{2A_n^+}\bm{\rho}_n^+=\frac{l_n}{2A_n^+}(\bm{r}-\bm{r}_+)&\bm{r}\in T_n^+\\
\frac{l_n}{2A_n^-}\bm{\rho}_n^-=\frac{l_n}{2A_n^-}(\bm{r}_--\bm{r})&\bm{r}\in T_n^-\\
0&\text{otherwise}
\end{cases}
\end{equation}

% figure
\begin{figure}[h]
	\includegraphics{pica.pdf}
	\caption{RWG 基函数几何参数示意图}
	\label{pica}
\end{figure}

% 2 figure
\begin{figure}[h]
\subfloat[]{
	\label{picb}
	\includegraphics[width=7.3cm]{picb.pdf}
}
\subfloat[]{
	\label{picc}
	\includegraphics[width=6.41cm]{picc.pdf}
}
\caption{调制高斯脉冲时域与频率波形，时域阻抗元素的存储技术也是时间步进算法并行化的关键技术之一。(a)调制高斯脉冲信号的时域波形；(b)调制高斯脉冲信号的频域波形}
\label{fig1}
\end{figure}

% veriloga
\begin{veriloga}[htbp]
\caption{Ideal Comparator for 14-bit SAR ADC}
\begin{lstlisting}
analog begin
    @(cross(V(in_p, in_n), 0))
        out = (V(in_p) > V(in_n));
end
\end{lstlisting}
\end{veriloga}

% footnote
\footnote{脚注序号“\ding{172}，……，\ding{180}”的字体是“正文”，不是“上标”，序号与脚注内容文字之间空1个半角字符，脚注的段落格式为：单倍行距，段前空0磅，段后空0磅，悬挂缩进1.5字符；中文用宋体，字号为小五号，英文和数字用Times New Roman字体，字号为9磅；中英文混排时，所有标点符号（例如逗号“，”、括号“（）”等）一律使用中文输入状态下的标点符号，但小数点采用英文状态下的样式“.”。}
（例如：FEKO、Ansys 等）的核心算法。由文献\cite{feng997he,clerc2010discrete,xiao2012yi}可知

